\section{Temporal Evaluation Research on a Separate Branch}
\label{sec:temporal-wip}

\subsection{Why chronological capture splits were not enough}

The master split already moves forward in time and keeps each capture intact.
This is much safer than randomly mixing ten-minute windows, but it gives only one
answer for one choice of fitting, calibration, and development dates. Device
behaviour can change with its operating state, daily routine, or traffic volume,
so a detector that works after one boundary may perform differently when that
boundary moves by a day.

The windows are also not independent repetitions. A cloud upload, scanning
burst, or attack can continue across several consecutive windows. Autocorrelation
measures this relationship; at a lag of one it compares each ten-minute window
with the window immediately after it. Within individual captures, Chromecast
total packet counts had an autocorrelation of about $0.86$, while Samsung
protocol proportions ranged from about $0.46$ to $0.60$. The anomaly scores were
also correlated across adjacent windows, at about $0.39$ for Chromecast and
$0.42$ for Samsung. These values do not prove that the Dirichlet--multinomial
model is wrong, but they do show that 144 windows from one day are not equivalent
to 144 independent days of evidence.

For the same reason, six adjacent attack windows may represent one continuing
attack rather than six separate incidents. Window-level precision and recall
remain useful, but the evaluation must also group consecutive alerts and labels
into episodes. Episode recall answers whether an attack was detected at all;
false-alert episodes per day describe how often an operator would be interrupted;
and detection delay shows how long the detector took to react. Results should
also remain visible for each capture so one poor day cannot be hidden inside a
large pooled confusion matrix.

The improved architecture therefore keeps complete captures together and repeats
the evaluation at several points in time. Each fold uses earlier normal captures
to fit the model, later normal captures to calibrate its threshold, and a
still-later capture to measure performance. The fitting history expands in the
next fold, while the final-test captures remain untouched.

Before these folds are used to select final settings, the temporal branch must be
combined with the current detector. The master branch should check the actual
time ranges inside every capture and record how overlapping or complementary
fragments were handled. Each fold should keep its own model, threshold, events,
and report, and results should be shown by capture and attack episode. If the
thesis reports confidence intervals, it should resample complete captures rather
than individual windows.

One D-Link data issue also remains: the October 8 and 9 fitting captures overlap
and must be resolved before the final model is fitted.

\subsection{Changes implemented on the temporal branch}

The \texttt{feature/temporal-model} branch orders every processed window by its
actual UTC start and end time. It stops when two captures overlap unless the
configuration explains how to resolve them, and it checks that fitting,
calibration, development, and final data occur in that order. For the known
exceptions, it combines three complementary Chromecast fragments. The report
records each decision.

The branch also defines expanding three-block folds:

\begin{center}
earlier fit history $\longrightarrow$ later threshold calibration
$\longrightarrow$ still-later validation
\end{center}

Each fold writes its model, threshold, events, and report to separate paths.
Besides the usual totals for individual windows, its report shows results for
each capture, groups consecutive attack windows into episodes, counts false-alert
episodes per day, and measures detection delay. Detector settings can then be
compared under a maximum false-alert rate chosen before the experiment. Final
captures remain outside every fold.

\subsection{Results from the temporal folds}

\begin{table}[H]
\centering
\small
\begin{tabularx}{\textwidth}{>{\raggedright\arraybackslash}p{0.18\textwidth}
>{\raggedright\arraybackslash}p{0.13\textwidth} r r r X}
\toprule
\textbf{Dataset} & \textbf{Fold} & \textbf{Windows} & \textbf{Window FP}
& \textbf{Episode recall} & \textbf{\shortstack[l]{False-alert\\episodes/day}} \\
\midrule
Chromecast & Benign 1 & 288 & 12 & -- & 4.00 \\
Chromecast & Benign 2 & 288 & 2 & -- & 1.00 \\
Chromecast & Labelled 3 & 474 & 4 & 0.70 & 0.91 \\
Samsung camera & Benign 1 & 81 & 1 & -- & 1.78 \\
Samsung camera & Labelled 2 & 556 & 4 & 0.82 & 0.78 \\
\bottomrule
\end{tabularx}
\caption{Committed temporal-fold reports. Episode recall is undefined in folds
without attack episodes.}
\label{tab:temporal-folds}
\end{table}

The two benign Chromecast periods contain the same number of windows, yet the
one-percent threshold produces 12 false alerts in one and two in the other. The
future false-alert rate is therefore not fixed by the calibration quantile. In
the labelled periods, the detector finds seven of ten Chromecast attack episodes
and 23 of 28 Samsung episodes. These figures describe the temporal branch only;
they cannot be compared directly with the adaptive results because that branch
was created before adaptation was implemented.

\subsection{The exact fold and timeline boundaries}

The three Chromecast folds expand fit history from October 10--13 to 10--15
and then 10--17, using later two-day calibration blocks and still-later
validation. Samsung's two folds move from May 28--29 fit/May 30 calibration/
May 31 validation to May 28--30 fit/May 31 calibration/June 1--4 validation.
They are complete-capture blocks, not random resamplings of adjacent windows.

The D-Link configuration on this branch excludes the overlapping October 8
capture. Its retained timeline validates eleven captures and 1,501 windows
without an unresolved overlap. Chromecast's declared complementary fragments
are resolved by explicit policy rather than silently counted twice.
These decisions change data provenance; they are not already present merely
because the reference configuration checks nominal dates.

The fold runner validates the actual timeline, fits and calibrates into
isolated paths, scores the fold's validation data, and writes capture/episode
reports. A selector groups candidates by score type, calibration quantile,
and window length, then applies a predeclared false-alert-episode budget before
ranking eligible attack-episode recall. Its source and reports are available
on the branch, but there is no claim of a new joint static/legacy/conservative
fold study in the current reference.

The reference branch and temporal branch also have different episode
capabilities. Episode counts, false-alert episodes/day, and detection delay
must not be attributed to the current reference evaluator, which still reports
window-level metrics only.
